\documentclass[10pt,journal,compsoc]{IEEEtran}
\usepackage{amsmath}
\usepackage{amsfonts}
\usepackage{amssymb}
\usepackage{graphicx}
\usepackage{url}

\begin{document}

\title{The Iscar Architecture: A Fully Decentralized, Off-Grid Resilient Matrix Fragmentation Network utilizing Zidon-Gideon Algebraic Anomalies}

\author{System Architect, Cryptographic Research Initiative}

\IEEEtitleabstractindextext{
\begin{abstract}
This paper introduces the Iscar Architecture, a novel cryptographic framework that eliminates centralized server dependencies and metadata trail exposure. By mapping message payloads into low-level machine code via a dynamically mutating multivariate finite vector space, the system provides high-grade resilience against asymmetric quantum cryptanalysis. We present the math behind a 15 by 15 parallel star proxy fragmentation matrix alongside local ad-hoc frequency handoffs, proving that total ephemeral communication is achievable across volatile mobile node constraints without persistent local or cloud logging footprint footprints.
\end{abstract}
\begin{IEEEkeywords}
Decentralized Mesh Networks, Asymmetric Cryptography, Off-Grid Communication, Quantum-Resistant Primitives.
\end{IEEEkeywords}}

\maketitle
\IEEEdisplaynontitleabstractindextext
\IEEEpeerreviewmaketitle

\section{Introduction}
Modern end-to-end encrypted communication applications remain vulnerable to systemic metadata extraction, state-sponsored routing infrastructure tap points, and cloud backup coercion vectors. The Iscar Architecture completely decouples data transit from static servers by leveraging a crowdsourced global network of smartphone processing units operating entirely within volatile execution environments.

\section{The Zidon-Gideon Cryptosystem Math}
Let the initial unencrypted string payload compile directly to raw machine logic binary sub-arrays denoted as $M_0$. When the data frames traverse across independent edge computing node matrices, each intervening device node $k$ forces a localized algorithmic mutation. The mathematical operation for each sequential encapsulation layer is formalized below:

\begin{equation}
Z_k(\xi) = [M_{k-1} \times 2^{\phi(k)}] \oplus \Gamma_{\text{local}} \pmod{\Lambda}
\end{equation}

Where the dynamic complexity multiplier parameter scales exponentially at variable interfaces:
\begin{equation}
\phi(k) \in \{25, 50, 200, \dots, 2^x\}
\end{equation}

\Lagrange Polynomial Reassembly Framework
The manual formalizes the recovery mechanics showing exactly how any threshold of m \ge 15 packets out of the total n = 30 split paths can synthesize the original payload perfectly using polynomial interpolation:

  $  \sum_{j=1}^{m} \prod_{m \neq j} \left[ \frac{x - x_m}{x_j - x_m} \right] \equiv M_0 $

The value $\Lambda$ represents a dynamic, multi-prime space generated distinctly for each standalone communication lifecycle channel. Because the terminal verification vector $\Gamma_{\text{local}}$ is derived purely from real-time physical micro-processor hardware clock entropy, the mathematical overhead prevents brute-force interception analysis over public routing arrays.

\section{15 $\times$ 15 Parallel Star Fragmentation Topology}
To completely strip metadata trails, payloads are shattered using non-linear threshold polynomial equations into 15 independent, mathematically disjoint data lines.

\begin{equation}
\sum_{j=1}^{m} \prod_{m \neq j} \left[ \frac{x - x_m}{x_j - x_m} \right] \equiv M_0 \quad (\text{for } m \ge 15)
\end{equation}

The architecture guarantees complete synthesis functionality at the target destination chip provided that any $m \ge 15$ of the 30 combined primary and mirror path segments reach the receiving device terminal safely.

\section{Conclusion and Future Review}
The Iscar network framework presents an un-killable, self-healing communication architecture. By relying exclusively on isolated RAM cells and local ad-hoc handoffs (Bluetooth Low Energy and Wi-Fi Direct), it mitigates cellular blackouts and hardware clean-room forensic extraction.

\bibliographystyle{IEEEtran}
\end{document}
